% STATUS: written 2026-11-05 (Lee 2001 read in full from the text layer;
% excerpt in excerpts/LeeTK2001SigmaCommuting.txt. Beidar--Chebotar and
% Posner 1957 are named through Lee 2001.)
\section{The $\sigma$-commuting form}\label{sec:functional-identities:centralizers-multipliers}

The centralizing maps of
Section~\ref{sec:functional-identities:commuting-maps} have a twisted
version, where the map is asked to commute with an endomorphism instead of
with the identity. Lee carried the characterization of
Section~\ref{sec:functional-identities:commuting-maps} over to this
setting and recorded it for semiprime rings
\cite[p.~2946]{LeeTK2001SigmaCommuting}. Throughout, $R$ is a semiprime
ring and $\sigma$ is an endomorphism of $R$
\cite[p.~2946]{LeeTK2001SigmaCommuting}. As in
Section~\ref{sec:functional-identities:commuting-maps}, $Z(R)$ is the
center, $C$ the extended centroid, $U$ the right Utumi quotient ring and
$Q$ the two-sided Utumi quotient ring of $R$
\cite[p.~2946]{LeeTK2001SigmaCommuting}.

A map $f$ of a subring $S$ of $R$ into $U$ is \emph{$\sigma$-commuting}
\cite[p.~2946]{LeeTK2001SigmaCommuting} when $[f(x),\sigma(x)]=0$ for all
$x$ in $S$. It is \emph{additive modulo $C$}
\cite[p.~2946]{LeeTK2001SigmaCommuting} when $f(x+y)-f(x)-f(y)$ lies in
$C$ for all $x,y$ in $S$. A right ideal of $R$ with zero left annihilator
is called \emph{large} \cite[p.~2946]{LeeTK2001SigmaCommuting}. The main
result generalizes the commuting-form Theorem~\ref{thm:fi:leetheorem2} to
the twisted case.

\begin{theorem}[{\cite[Theorem~2, p.~2948]{LeeTK2001SigmaCommuting}}]\label{thm:fi:leesigmamain}
Let $R$ be a semiprime ring, $I$ a subring of $R$ and $\sigma$ an
endomorphism of the ring $R$ such that $\sigma(I)$ contains a large ideal
of $R$. Then every $\sigma$-commuting map $f\colon I\to U$, which is
additive modulo $C$, must be of the form $f(x)=\lambda\sigma(x)+\zeta(x)$
for all $x\in I$, where $\lambda\in C$ and $\zeta\colon I\to C$.
\end{theorem}

Under the weaker hypothesis that $\sigma(I)$ contains only a large right
ideal, the same form holds with values in the two-sided quotient $Q$
\cite[Theorem~1, p.~2947]{LeeTK2001SigmaCommuting}. When $\sigma$ is an
epimorphism the hypothesis is automatic on the whole ring, so every
$\sigma$-commuting map of $R$ that is additive modulo $C$ has the form
$f(x)=\lambda\sigma(x)+\zeta(x)$
\cite[Corollary~1, p.~2948]{LeeTK2001SigmaCommuting}. The proof of
Theorem~\ref{thm:fi:leesigmamain} rests on a biderivation identity,
Lemma~\ref{thm:fi:leesigmalemma}, that is itself a functional identity.

\begin{lemma}[{\cite[Lemma~1, p.~2947]{LeeTK2001SigmaCommuting}}]\label{thm:fi:leesigmalemma}
Let $R$ be a semiprime ring with an endomorphism $\sigma$ and let
$D\colon S\times S\to U$ be a $(\sigma,\sigma)$-biderivation, where $S$ is
a subring of $R$. Then $[\sigma(x),\sigma(y)]D(u,v)=D(x,y)[\sigma(u),\sigma(v)]$
for all $x,y,u,v$ in $S$.
\end{lemma}

The $\sigma$-commuting form forces a Posner-type conclusion for twisted
derivations \cite[Theorem~3, p.~2948]{LeeTK2001SigmaCommuting}. Let
$\sigma$ and $\tau$ be endomorphisms of $R$; a
$(\sigma,\tau)$-derivation is an additive map $d$ with
$d(xy)=d(x)\sigma(y)+\tau(x)d(y)$
\cite[p.~2948]{LeeTK2001SigmaCommuting}.

\begin{theorem}[{\cite[Theorem~3, p.~2948]{LeeTK2001SigmaCommuting}}]\label{thm:fi:leesigmat3}
Let $R$ be a semiprime ring with $d$ a $(\sigma,\tau)$-derivation of $R$.
Suppose that $I$ is a subring (right ideal) of $R$ such that both
$\sigma(I)$ and $\tau(I)$ contain some large right ideals of $R$. Suppose
that $[d(x),\sigma(x)]=0$ for all $x\in I$. Then $d(I)$, and likewise $d(R)$, is
contained in a central ideal of $R$.
\end{theorem}

For a prime ring the conclusion of Theorem~\ref{thm:fi:leesigmat3}
becomes commutativity, since a nonzero $(\sigma,\tau)$-derivation with
$[d(x),\sigma(x)]=0$ on a right ideal forces $R$ to be commutative
\cite[Corollary~4, p.~2950]{LeeTK2001SigmaCommuting}. This is the
semiprime and prime form of the theorem of Posner for twisted derivations, and
it recovers the classical statement of 1957 on specializing
$\sigma=\tau=\mathrm{id}$ \cite[p.~2946]{LeeTK2001SigmaCommuting}. On a
noncentral Lie ideal the $\sigma$-centralizing form still holds, with two
low-rank exceptions, as Theorem~\ref{thm:fi:leesigmat4} states.

\begin{theorem}[{\cite[Theorem~4, p.~2950]{LeeTK2001SigmaCommuting}}]\label{thm:fi:leesigmat4}
Let $R$ be a prime ring, $L$ a noncentral Lie ideal of $R$ and $\sigma$ an
epimorphism of $R$ with $\sigma(L)$ not contained in $Z(R)$. Then every
$\sigma$-centralizing map $f\colon L\to U$, which is additive modulo $C$,
must be of the form $f(x)=\lambda\sigma(x)+\zeta(x)$ for all $x\in L$,
where $\lambda\in C$ and $\zeta\colon L\to C$ unless $\dim_C RC=4$ or $9$.
\end{theorem}

Lee quotes Theorem~\ref{thm:fi:leesigmat4} as an immediate consequence of
a general functional-identities result of Beidar and Chebotar
\cite[p.~2950]{LeeTK2001SigmaCommuting}. This last statement is the point
where the twisted commuting theory meets the $d$-free method of
Section~\ref{sec:functional-identities:d-free}
\cite[p.~2950]{LeeTK2001SigmaCommuting}.

For a nonzero generalized skew derivation $F$ on a noncentral Lie ideal of a ring, the $n$-commuting condition $[F(x^n),x]=0$ is shown to force $F$ to vanish on the ideal under stated hypotheses~\cite{WlPrimeAndSemiprime2018}.

An $n$-centralizing generalized derivation on a semiprime ring is shown to behave as the commuting case does, with an application to $C^*$-algebras~\cite{WlOnN-centralizingGeneralize2012}.

An $n$-centralizing generalized derivation on a semiprime ring is shown to behave as the commuting case does, with an application to $C^*$-algebras~\cite{WlOnN-centralizingGeneralize2012}.

Identities involving two generalized derivations on a noncentral Lie ideal of a prime ring are shown to force the derivations into a standard relation~\cite{WlIdentitiesWithGeneralized2012}.
